% ============ 结论 ============
\chapter*{结\quad 论}
\addcontentsline{toc}{chapter}{结\quad 论}
\markboth{结论}{结论}

本课题围绕智能眼镜普及背景下显示器敏感信息被偷拍识别的安全威胁，设计并实现了一个基于视觉暂留与时间分片掩模的显示器隐私保护概念验证系统。系统利用人眼约 50ms 视觉积分窗口与相机高速快门采样的根本差异，将每帧图像在时域拆分为随机点阵子帧并叠加时域互补的对抗噪声，使人眼积分后感知完整画面，而相机任一瞬时采样只能捕获碎片化残片。

\noindent\textbf{（一）主要工作与创新点}

本课题的主要工作与创新点包括：

（1）实现了由 ChaCha20 密码学安全伪随机数生成器与 HKDF 周期密钥派生驱动的随机点阵掩模生成算法，通过拒绝采样消除取模偏置、通过卡方检验保证分布均匀，使 $n$ 个子帧严格满足完备性与互斥性约束、周期间统计独立。

（2）实现了时域互补的对抗噪声注入算法，将 FGSM/PGD 基础噪声精确分解为 $\sum_k N_k = 0$ 的互补子噪声，使噪声在人眼积分后完全抵消、在单帧内最大化识别错误率；实验进一步证明对抗噪声是单帧防御抵抗 inpainting 攻击的关键支撑。

（3）实现了包含伪随机时序置换、时序令牌、反色帧/黑帧防御的时序控制，以及 GPU 优先、软件回退的渲染链路，并配套 HDR 亮度补偿、多显示器同步、视觉疲劳策略等辅助模块，以 166 项单元测试保证算法正确性。

（4）建立了多引擎 OCR 与目标检测攻防量化评测体系，证明系统在 $n=4$、240Hz 下视觉无感（FPI 0.030、$\Delta E\approx0.37$）；在 120 样本文本语料上，Tesseract、EasyOCR、PaddleOCR 的单子帧 OCR 准确率均约为 0.0\%，配对降幅达到 93.9\%--94.9\%，词级准确率、exact-match 与敏感 token 恢复率均为 0.0\%，且对目标威胁模型内的抽帧、流式、卷帘快门、插反色帧长曝光等攻击保持有效。

（5）诚实揭示并以实验证明了线性时域叠加方案的原理性局限——覆盖完整周期的时域平均、相位搜索与单通道最大值重构均可还原图像；在 120 样本文本语料上，逐样本最强攻击的字符恢复率达到 95.2\%、exact-match 达到 91.7\%、敏感 token 恢复率达到 98.6\%，明确界定了系统"防被动抽帧/流式、不防主动完整周期重构"的防护边界，这一结论对同类方案具有普遍的参考价值。

\noindent\textbf{（二）不足与展望}

本课题仍存在以下不足，有待进一步研究：

（1）本系统是应用/算法层的概念验证原型，技术方案中操作系统合成器层、GPU 驱动层的透明像素注入尚未实现；下一步可基于 DXGI、KMS/DRM、CoreDisplay 等接口实现驱动级注入，达到对全屏应用的透明覆盖与生产级性能。

（2）受实验条件限制，本课题在普通 60--120Hz 显示器上以缩小的拆分因子演示并外推高刷指标，未在 240Hz/480Hz 真机上开展受试者主观实验；下一步可补充真机验证与人因评测。

（3）针对完整周期对齐平均这一根本局限，下一步可探索利用物理 LCD 视角响应的不可逆视角差异化、引入人眼与相机积分差异的非线性时域调制、随机化周期长度与相位等缓解手段，突破纯线性方案的能力边界。

综上，本课题以纯软件方式验证了时间分片掩模显示隐私保护的可行性与有效边界，对智能眼镜防护、机密信息终端等场景具有应用价值，并为后续驱动层产品化与专利申请提供了可靠的实验依据。
